
%\needspace{4\baselineskip}
\color{uniblue}{\subsection{Логическая защита шин (ЛЗШ)}}
\color{black}

\begin{enumerate}

\item
Логическая защита шин предназначена для отключения вводного и секционного выключателей при коротких замыканиях на секции шин с минимальной выдержкой времени срабатывания. 

\item
В устройстве логическая защита шин представлена функцией «ЛЗШ». 

\item
В  качестве  ИО  ЛЗШ  применяется  ИО  тока  максимального  действия,  принцип  работы  которого заключается в пофазном сравнении действующего значения основной гармоники тока с уставкой срабатывания и коэффициентом возврата порогового равным 0,95. Характеристика срабатывания измерительного органа –- независимая.

\item
В  случае  пуска  МТЗ  отходящих  присоединений  первой или второй секции шин действие функции блокируется. При  условии  срабатывания  любого  фазного  ИО  ЛЗШ  и  отсутствии  пуска  МТЗ  отходящих присоединений первой или второй секции шин происходит пуск ЛЗШ. Срабатывание  функции  и  отключение выключателя происходит по истечении независимой выдержки времени таймера Т1. 

\item
В функции реализовано две схемы работы алгоритма: последовательная и параллельная. Они учитывают контакты выходных реле отходящих линий первой или второй секции шин, которые сигнализируют о пуске  МТЗ. Выбор схемы подключения осуществляется с помощью  уставки  SGF4. Поясняющие схемы работы ЛЗШ представлены на рисунках \ref{lzsh:appimg1}, \ref{lzsh:appimg2} приложения \ref{app:lzsh}. 

\item
Последовательная  схема  предусматривает  использование  нормально  замкнутых  контактов  реле «Пуск  МТЗ» от устройств защиты отходящих присоединений, параллельная схема предусматривает использование нормально разомкнутых контактов реле.

\item
При длительном присутствии сигналов на входах «Блокировка ЛЗШ», по истечении нерегулируемой выдержки времени таймера Т2 (5 с) активируется выходной сигнал «ЛЗШ: Неисправность», который указывает на наличие неисправности в схеме ЛЗШ.

\item
Ввод функции в работу на этапе параметрирования устройства осуществляется программным переключателем SGF1 «Ввод\_функции». Информация о введенной функции отображается выходным сигналом «ЛЗШ: Ввод».

\item\sloppy
Функция может быть выведена из работы оперативно путем активации сигнала «ОВ ЛЗШ». Выведенное состояние функции характеризуется активным состоянием сигнала «ЛЗШ: Оперативный вывод».

\item
При активном состоянии входного сигнала «ЛЗШ на сигн» (ключ «ЛЗШ на сигн» введен) действие функции переводится на сигнал.

\item
Алгоритм работы функции выполнен в соответствии с рисунком \ref{lzsh:img1}. В таблице \ref{lzsh:tbl1} приведены параметры, необходимые для настройки функции.

\vspace{1mm}
\begin{figure}[H]
\centering
\includegraphics[width=1\textwidth,height=1\textheight,keepaspectratio]{\fbpath/302. ЛЗШ/_latex/img/lzsh.pdf}
\captionof{figure}{Функциональная схема ЛЗШ}\label{lzsh:img1}
\end{figure}

\end{enumerate}
